% Unidad U013. Biografía estructural completa del carácter exponencial.

\label{gfp:b:chap:biografia-completa-e}

\section{Finite orbital selector: existence and uniqueness of compatible emission}

Let \(\mathcal W=\mathbb F_3^6\). For an emission
\(U=(u_1,\ldots,u_6)\), the visible word is
\[
 \Psi(U)=((-u_1)\bmod3,\ldots,(-u_6)\bmod3).
\]
On \(\mathcal W\), the dihedral action is given by
\[
\begin{aligned}
 r(u_1,u_2,u_3,u_4,u_5,u_6)
 &=(u_3,u_1,u_2,u_5,u_6,u_4),\\
 s(u_1,u_2,u_3,u_4,u_5,u_6)
 &=(u_4,u_5,u_6,u_1,u_2,u_3).
\end{aligned}
\]
For \(w\in\mathcal W\), define
\[
 n(w)=\#\Psi^{-1}(w),
 \qquad
 M(w)=\sum_{U\in\Psi^{-1}(w)}\operatorname{mult}(U),
\]
and
\[
 \nu(w)=
 \bigl(
 \#\{j:w_j=0\},
 \#\{j:w_j=1\},
 \#\{j:w_j=2\}
 \bigr).
\]

\begin{theorem}[Structural selection of the exponential channel]
\label{gfp:b:thm:seleccion-estructural-e}
Among the \(43\) orbits of the visible catalog, there is a unique orbit whose
fibers simultaneously satisfy:
\[
 n(w)=1,\qquad M(w)=144,\qquad
 \mathfrak D(w)=\{0,3,6\},\qquad
 \nu(w)=(2,3,1).
\]
The gauge
\[
 \operatorname{diag}_c=6,
 \qquad \operatorname{card}=ES
\]
select in it
\[
 \boxed{w_6^e=201101.}
\]
\end{theorem}

\begin{proof}
The four conditions are evaluated over the \(243\) words in the catalog, not over a sample. A single orbit satisfies the joint predicate. In that orbit, exactly one emission has both diagonal \(6\) and orientation \(ES\). The enumeration uses only fields from the APP--TRIT catalog and does not consult an expansion of \(e\).
\end{proof}

\section{Block elevation and successive boundary compatibility}

The recurrence \(L_0,L_0,L_1,L_1\) produces
\[
 w_{30}^{e}
 =201101|121221|102011|012222|102011,
\]
and the joint boundary contributes
\[
 u_5^e=021222.
\]
The record of thirty-six trits remains
\[
 w_{36}^{e}
 =201101|121221|102011|012222|102011|021222.
\]

\section{Decimal positional signature and compatibility square}

The initial reading of twelve digits is
\[
 \mathbf d_e^{(12)}
 =718281828459
 =7\,\underbrace{1828\,1828}_{\text{local square}}\,459.
\]
The central equality is the concatenation
\[
 1828\mathbin{\|}1828,
\]
not an arithmetic product nor the supposed beginning of a period.

\begin{definition}[Positional square]
A word \(d_1\cdots d_{12}\) contains a square of length \(\ell\)
at position \(p\) when
\[
 d_p\cdots d_{p+\ell-1}
 =d_{p+\ell}\cdots d_{p+2\ell-1}.
\]
The position and the lateral contexts belong to the signature.
\end{definition}

\begin{table}[htbp]
\centering
\caption{Comprehensive census of squares in twelve-digit readings.}
\label{gfp:b:tab:censo-cuadrados-e}
\begin{tabular}{c*{6}{r}}
\toprule
Length \(\ell\)&1&2&3&4&5&6\\
\midrule
Occurrences&240&21&3&2&0&0\\
Words that contain them&153&19&3&2&0&0\\
\bottomrule
\end{tabular}
\end{table}

The other word with a square of length four is
\[
 575157518472
 =\underbrace{5751\,5751}_{\text{initial position }1}\,8472.
\]

\begin{proposition}[Positional uniqueness]
\label{gfp:b:prop:cuadrado-posicional-e}
The reading of channel \(e\) is the only one that has a square of length
four after a prefix of length one and with contexts \(7\) and \(459\).
\end{proposition}

\begin{proof}
The census leaves only two words with squares of length four. One starts
it at the first position; the other after the prefix \(7\). The
position and contexts separate the two possibilities.
\end{proof}

\section{Phase return with a change of state and conservation of traversed memory}

In the elevation
\[
 201101|121221|102011|012222|102011
\]
the third and the fifth blocks coincide:
\[
 B_3=B_5=102011.
\]
Their prestates modulo \(27\) are, however,
\[
 p_3=(25,11,7,17,8,16),
 \qquad
 p_5=(7,8,4,23,26,7),
\]
and the subsequent quotients:
\[
 q_3=(6,13,9,2,13,1),
 \qquad
 q_5=(15,0,3,19,23,25).
\]
\(p_3\ne p_5\) and \(q_3\ne q_5\) hold.

\begin{figure}[htbp]
\centering
\begin{tikzpicture}[
 node distance=11mm and 7mm,
 every node/.style={font=\small},
 block/.style={draw,rounded corners,minimum width=17mm,minimum height=8mm},
 state/.style={draw, dashed, rounded corners, align=center, inner sep=3pt},
 >=stealth
]
\node[block] (b1) {\(201101\)};
\node[block,right=of b1] (b2) {\(121221\)};
\node[block,right=of b2] (b3) {\(102011\)};
\node[block,right=of b3] (b4) {\(012222\)};
\node[block,right=of b4] (b5) {\(102011\)};
\draw[->] (b1)--(b2);
\draw[->] (b2)--(b3);
\draw[->] (b3)--(b4);
\draw[->] (b4)--(b5);
\node[state,below=of b3] (p3)
 {\(p_3=(25,11,7,17,8,16)\)\\
  \(q_3=(6,13,9,2,13,1)\)};
\node[state,below=of b5] (p5)
 {\(p_5=(7,8,4,23,26,7)\)\\
  \(q_5=(15,0,3,19,23,25)\)};
\draw[->] (p3)--(b3);
\draw[->] (p5)--(b5);
\draw[<->,thick] (b3) to[bend left=28]
 node[above]{same word visible} (b5);
\end{tikzpicture}
\caption{Visible return of \(102011\) with renewal of the prestate and the
quotient.}
\label{gfp:b:fig:ruptura-memoria-e}
\end{figure}

\begin{theorem}[Visible-phase return without periodicity of the state]
\label{gfp:b:thm:retorno-visible-no-periodico-e}
\(B_3=B_5\) is a return of the visible projection, not a periodic orbit
of the enriched state.
\end{theorem}

\begin{proof}
The return of the complete state would require equality of all its coordinates.
The prestate and quotient inequalities exclude it. Nor does
\(1828|1828\) begin a decimal period: after the second block,
\(4\) appears, whereas a periodic continuation would require \(1\).
\end{proof}

\section{Subsequent factorial certification of the coinductive exponential section}

Let \(\mathcal A\) be a unital commutative algebra over \(\mathbb Q\), and
let \(P_t\in\mathcal A[[t]]\) be the formal propagator. Temporal
concatenation imposes
\[
 P_{s+t}=P_sP_t,
 \qquad
 P_0=1.
\]
We write
\[
 P_t=\sum_{n=0}^{\infty}a_nt^n.
\]
The unit fixes \(a_0=1\), and normalized elementary transport fixes
\(a_1=1\).

\begin{theorem}[Factorial recurrence]
\label{gfp:b:thm:recurrencia-factorial-e}
\[
 (n+1)a_{n+1}=a_n
 \qquad(n\geq0),
\]
and, therefore,
\[
 a_n=\frac1{n!},
 \qquad
 P_t=\sum_{n=0}^{\infty}\frac{t^n}{n!}.
\]
En \(t=1\),
\[
 \boxed{P_1=e.}
\]
\end{theorem}

\begin{proof}
The coefficient of \(s^nt\) in \(P_{s+t}\) is
\[
 \binom{n+1}{n}a_{n+1}=(n+1)a_{n+1}.
\]
In \(P_sP_t\) it is \(a_na_1=a_n\). Equality of series produces the
recurrence. Induction from \(a_0=1\) gives \(a_n=1/n!\).
\end{proof}

\begin{corollary}[Differential equation]
\[
 \frac{d}{dt}P_t=P_t,
 \qquad P_0=1.
\]
\end{corollary}

\begin{proof}
\[
 P_t'=\sum_{n\geq0}(n+1)a_{n+1}t^n
 =\sum_{n\geq0}a_nt^n=P_t.
\]
\end{proof}

\section{Cofinal Forking of Rational Tangles and Convergence to the Published Value}

Let
\[
 S_N=\sum_{n=0}^{N}\frac1{n!}.
\]

\begin{theorem}[Explicit factorial bound]
\label{gfp:b:thm:cota-factorial-e}
For \(N\geq1\),
\begin{equation}
 \boxed{
 S_N<e<
 S_N+\frac{N+2}{(N+1)(N+1)!}.}
 \label{gfp:b:eq:horquilla-factorial-e}
\end{equation}
\end{theorem}

\begin{proof}
The rest is
\[
 R_N=\sum_{k=0}^{\infty}\frac1{(N+1+k)!}>0.
\]
Para \(k\geq0\),
\[
 (N+1+k)!
 \geq(N+1)!(N+2)^k.
\]
Therefore,
\[
 R_N
 \leq\frac1{(N+1)!}
 \sum_{k=0}^{\infty}\frac1{(N+2)^k}
 =\frac{N+2}{(N+1)(N+1)!}.
\]
The inequality is strict from \(k=2\) onward.
\end{proof}

We define
\[
 J_N^{(e)}=
 \left[
 S_N,\,
 S_N+\frac{N+2}{(N+1)(N+1)!}
 \right]
\]
and, for the fractional part,
\[
 \widetilde J_N^{(e)}=J_N^{(e)}-2.
\]
The width converges to zero. For each \(M\), there exists \(N\) such that
\[
 \operatorname{diam}J_N^{(e)}<10^{-M}.
\]
If the endpoints belong to the same decimal cell, the first \(M\)
digits are certified by rational arithmetic.

\Needspace{30\baselineskip}
\section{Coinductive prolongation of the exponential section through ternary subcylinders of the Topological Phase Kernel}

At level \(K\), TPK and the relation \(\operatorname{Ext}\) have already generated
a unique compatible subcylinder. The smallest \(N=N_e(K)\) for which
the factorial bracket is contained in that subcylinder is taken. \(N_e(K)\) is an
order of analytic certification, not a generating index or a rule for
selecting the fiber.

The TPK cylinders are compatible under truncation independently of the
semigroup. The semigroup theorem subsequently identifies their Archimedean
evaluation with \(e-2\); the integer part recovers \(e\). If the orbit,
\(\operatorname{Ext}\) or truncation fails, generation fails; if the
recurrence or the factorial bound fails, this analytic certificate fails, not the
APP--TRIT--TPK genealogy.

\section{Internal certificate of orbital uniqueness, factorial recurrence, convergence, and cylindrical prolongation}

\begin{remark}[Certified and validity control: Falsification of the exponential channel]
The construction is refuted if:
\begin{enumerate}
 \item another orbit satisfies the complete finite predicate;
 \item the signature \(7[1828\,1828]459\) is not positionally unique;
 \item the prestates or quotients of the two blocks \(102011\) coincide;
 \item the semigroup does not force \((n+1)a_{n+1}=a_n\);
 \item the bound \eqref{gfp:b:eq:horquilla-factorial-e} fails;
 \item a level does not locate a compatible extension uniquely;
 \item a tabulated expansion intervenes before recognition.
\end{enumerate}
\end{remark}
